% ============ §5 The Hilbert Nullstellensatz and first steps in dimension theory (原书 pp.30-37) ============
\section{The Hilbert Nullstellensatz and first steps in dimension theory}\label{sec:5}\index{Hilbert Nullstellensatz, section 5}\index{dimension theory, section 5}

Let $X$ be a topological space; we consider strictly decreasing (or
strictly increasing) chains $Z_0 \supsetneq Z_1 \supsetneq \cdots \supsetneq Z_r$ of length $r$ of irreducible closed
subsets of $X$. The supremum of the lengths, taken over all such chains, is
called the \emph{combinatorial dimension}\index{combinatorial dimension} of $X$ and denoted $\dim X$. If $X$ is a
Noetherian space then there are no infinite strictly decreasing chains, but it
can nevertheless happen that $\dim X = \infty$.

Let $Y$ be a subspace of $X$. If $S \subset Y$ is an irreducible closed subset of $Y$ then
its closure in $X$ is an irreducible closed subset $S' \subset X$ such that $S' \cap Y = S$.
Indeed, if $S' = V \cup W$ with $V$ and $W$ closed in $X$ then $S = (V \cap Y) \cup (W \cap Y)$,
so that say $S = V \cap Y$, but then $V = S'$. It follows easily from this that
$\dim Y \le \dim X$.

Let $A$ be a ring. The supremum of the lengths $r$, taken over all strictly
decreasing chains $p_0 \supsetneq p_1 \supsetneq \cdots \supsetneq p_r$ of prime ideals of $A$, is called the \emph{Krull
dimension}\index{Krull dimension}, or simply the \emph{dimension} of $A$, and denoted $\dim A$. As one sees
easily from Ex.~\ref{ex:4.10-ex}, the Krull dimension of $A$ is the same thing as the
combinatorial dimension of $\Spec A$. For a prime ideal $p$ of $A$, the supremum
of the lengths, taken over all strictly decreasing chains of prime ideals
$p = p_0 \supsetneq p_1 \supsetneq \cdots \supsetneq p_r$ starting from $p$, is called the \emph{height} of $p$,\index{height of a prime ideal} and denoted
$\htop p$; (if $A$ is Noetherian it will be proved in Theorem~\ref{thm:13.5} that $\htop p < \infty$).
Moreover, the supremum of the lengths, taken over all strictly increasing
chains of prime ideals $p = p_0 \subsetneq p_1 \subsetneq \cdots \subsetneq p_r$ starting from $p$, is called the
\emph{coheight} of $p$, and written $\coht p$. It follows from the definitions that
\[
  \htop p = \dim A_p, \quad \coht p = \dim A/p \quad \text{and} \quad \htop p + \coht p \le \dim A.
\]

\begin{remark}
In more old-fashioned terminology $\htop p$ was usually called the
\emph{rank} of $p$, and $\coht p$ the \emph{dimension} of $p$; in addition, Nagata [N1] calls
$\dim A$ the \emph{altitude} of $A$.
\end{remark}

\begin{example}\label{ex:5.1}
The prime ideals in the ring $\Z$ of rational integers are the ideals
$p\Z$ generated by the primes $p = 2, 3, 5, \ldots$, together with $(0)$. Hence, every $p\Z$
is a maximal ideal, and $\dim \Z = 1$. More generally, any principal ideal
domain which is not a field is one-dimensional.
\end{example}

\begin{example}\label{ex:5.2}
An Artinian ring is zero-dimensional; indeed, we have seen in
the proof of Theorem~\ref{thm:3.2} that there are only a finite number of maximal
ideals $p_1, \ldots, p_n$, and that the product of all of these is nilpotent. If then $p$ is a
prime ideal, $p \supset (0) = (p_1 \cdots p_n)^\nu$ so that $p \supset p_i$ for some $i$; hence, $p = p_i$, so
that every prime ideal is maximal.
\end{example}

\begin{example}\label{ex:5.3}
A zero-dimensional integral domain is just a field.
\end{example}

\begin{example}\label{ex:5.4}
The polynomial ring $k[X_1, \ldots, X_n]$ over a field $k$ is an integral
domain, and since
\[
  k[X_1, \ldots, X_n]/(X_1, \ldots, X_i) \cong k[X_{i+1}, \ldots, X_n],
\]
$(X_1, \ldots, X_i)$ is a prime ideal of $k[X_1, \ldots, X_n]$. Thus
\[
  (0) \subset (X_1) \subset (X_1, X_2) \subset \cdots \subset (X_1, \ldots, X_n)
\]
is a chain of prime ideals of length $n$, and $\dim k[X_1, \ldots, X_n] \ge n$. In fact we
will shortly be proving that equality holds.
\end{example}

For an ideal $I$ of a ring $A$ we define the \emph{height of $I$} to be the infimum of the
heights of prime ideals containing $I$:
\[
  \htop I = \inf\{\htop p \mid I \subset p \in \Spec A\}.
\]
Here also we have the inequality
\[
  \htop I + \dim A/I \le \dim A.
\]
If $M$ is an $A$-module we define the \emph{dimension of $M$} by
\[
  \dim M = \dim(A/\ann(M)).
\]
If $M$ is finitely generated then $\dim M$ is the combinatorial dimension of the
closed subspace $\Supp(M) = V(\ann(M))$ of $\Spec A$.

A strictly increasing (or decreasing) chain $p_0, p_1, \ldots$ of prime ideals is said
to be \emph{saturated}\index{saturated chain} if there do not exist prime ideals strictly contained between
any two consecutive terms. We say that $A$ is a \emph{catenary ring}\index{catenary ring} if the following
condition is satisfied: for any prime ideals $p$ and $p'$ of $A$ with $p \subset p'$, there
exists a saturated chain of prime ideals starting from $p$ and ending at $p'$, and
all such chains have the same (finite) length.

If a local domain $(A, m)$ is catenary then for any prime ideal $p$ we have
$\htop p + \coht p = \dim A$. Conversely, if $A$ is a Noetherian local domain and
this equality holds for all $p$ then $A$ is catenary (Ratliff [3], 1972); the proof
of this is difficult, and we postpone it to Theorem 31.4. Practically all the
important Noetherian rings arising in applications are known to be
catenary; the first example of a non-catenary Noetherian ring was
discovered in 1956 by Nagata [15].

We now spend some time discussing the elementary theory of dimensions
of rings which are finitely generated over a field $k$.

\begin{theorem}\label{thm:5.1}
Let $k$ be a field, $L$ an algebraic extension of $k$ and
$a_1, \ldots, a_n \in L$; then
\begin{enumerate}[(i)]
  \item $k[a_1, \ldots, a_n] = k(a_1, \ldots, a_n)$.
  \item Write $\varphi\colon k[X_1, \ldots, X_n] \to k[a_1, \ldots, a_n]$ for the homomorphism
        over $k$ which maps $X_i$ to $a_i$; then $\Ker \varphi$ is the maximal ideal generated
        by $n$ elements of the form $f_1(X_1), f_2(X_1, X_2), \ldots, f_n(X_1, \ldots, X_n)$, where
        each $f_i$ can be taken to be monic in $X_i$.
\end{enumerate}
\end{theorem}

\begin{proof}
Let $f_1(X)$ be the minimal polynomial of $a_1$ over $k$; then $(f_1(X_1))$ is a
maximal ideal of $k[X_1]$, so that $k[a_1] \cong k[X_1]/(f_1(X_1))$ is a field, and hence
$k[a_1] = k(a_1)$. Now let $\varphi_2(X)$ be the minimal polynomial of $a_2$ over $k(a_1)$;
then since $k(a_1) = k[a_1]$, the coefficients of $\varphi_2$ can be expressed as
polynomials in $a_1$, and there is a polynomial $f_2 \in k[X_1, X_2]$, monic in $X_2$,
such that $\varphi_2(X_2) = f_2(a_1, X_2)$. Thus
\[
  k[a_1, a_2] = k[a_1][a_2] = k(a_1)[a_2] \cong k(a_1)[X_2]/(f_2(a_1, X_2)).
\]
Proceeding in the same way, for $1 \le i \le n$ there is an
$f_i(X_1, \ldots, X_i) \in k[X_1, \ldots, X_i]$, monic in $X_i$, such that
\[
  k[a_1, \ldots, a_i] = k(a_1, \ldots, a_i)
\]
\[
  \cong k(a_1, \ldots, a_{i-1})[X_i]/(f_i(a_1, \ldots, a_{i-1}, X_i)).
\]

Now if $P(X) \in k[X_1, \ldots, X_n]$ is in the kernel of $\varphi$, we have $\varphi(P) =
P(a_1, \ldots, a_n) = 0$, so that $P(a_1, \ldots, a_{n-1}, X_n)$ is divisible by $f_n(a_1, \ldots,
a_{n-1}, X_n)$; dividing $P(X_1, \ldots, X_n)$ as a polynomial in $X_n$ by the
monic polynomial $f_n(X_1, \ldots, X_n)$ and letting $R_n(X_1, \ldots, X_n)$ be the remain-
der, we can write $P = Q_n f_n + R_n$, with $R_n(a_1, \ldots, a_{n-1}, X_n) = 0$. Similarly,
dividing $R_n(X_1, \ldots, X_n)$ as a polynomial in $X_{n-1}$ by $f_{n-1}(X_1, \ldots, X_{n-1})$ and
letting $R_{n-1}(X_1, \ldots, X_n)$ be the remainder, we get
\[
  R_n = Q_{n-1}f_{n-1} + R_{n-1},
\]
with
\[
  R_{n-1}(a_1, \ldots, a_{n-2}, X_{n-1}, X_n) = 0.
\]
Proceeding in the same way we get $P = \sum Q_i f_i + R$, with $R(X_1, \ldots, X_n) = 0$;
that is $R = 0$ and $P = \sum Q_i f_i$, so that $\Ker \varphi = (f_1, f_2, \ldots, f_n)$.
\end{proof}

The following theorem can be regarded as a converse of Theorem 1, (i).

\begin{theorem}\label{thm:5.2}
Let $k$ be a field and $A = k[\alpha_1, \ldots, \alpha_n]$ an integral domain, and
write $r = \trdeg_k A$ for the transcendence degree of $A$ (that is, of its field of
fractions) over $k$. Then if $r > 0$, $A$ is not a field.
\end{theorem}

\begin{proof}
Suppose that $\alpha_1, \ldots, \alpha_r$ is a transcendence basis for $A$ over $k$, and set
$K = k(\alpha_1, \ldots, \alpha_r)$. Then since $\alpha_{r+1}, \ldots, \alpha_n$ are algebraic over $K$, there exist
polynomials $f_i(X_{r+1}, \ldots, X_i) \in K[X_{r+1}, \ldots, X_i]$, monic of degree $d_i$ in $X_i$, such that
\[
  k[\alpha_{r+1}, \ldots, \alpha_n] \cong K[X_{r+1}, \ldots, X_n]/(f_{r+1}, \ldots, f_n)
\]
and
\[
  d_i = [K(\alpha_{r+1}, \ldots, \alpha_i) : K(\alpha_{r+1}, \ldots, \alpha_{i-1})] \quad \text{for } i > r.
\]
The coefficients of $f_i$ are in $K$, so that for suitable $0 \neq g \in k[\alpha_1, \ldots, \alpha_r]$ we
have $gf_i \in k[\alpha_1, \ldots, \alpha_r][X_{r+1}, \ldots, X_n]$. In other words, if we set $B =
k[\alpha_1, \ldots, \alpha_r, g^{-1}]$ then the $f_i$ are polynomials with coefficients in $B$. We
are now going to show that $A[g^{-1}] = B[\alpha_{r+1}, \ldots, \alpha_n]$ is a free module over
$B$ with basis $\{\prod_{i>r} \alpha_i^{e_i} \mid 0 \le e_i < d_i\}$. Every element of $B[\alpha_{r+1}, \ldots, \alpha_n]$ can
be written as $P(\alpha_{r+1}, \ldots, \alpha_n)$ for some $P \in B[X_{r+1}, \ldots, X_n]$; dividing $P$ as a
polynomial in $X_n$ by $f_n$, and replacing $P$ by the remainder, we can assume
that $P$ has degree at most $d_n - 1$ in $X_n$; then dividing $P$ as a polynomial
in $X_{n-1}$ by $f_{n-1}$, and replacing by the remainder, we can assume that $P$
has degree at most $d_{n-1} - 1$ in $X_{n-1}$. Proceeding in the same way, we
can assume that $P$ has degree at most $d_i - 1$ in $X_i$ for each $i$; in addition,
the elements $\{\alpha_i^e \mid 0 \le e < d_i\}$ are linearly independent over $K(\alpha_{r+1}, \ldots, \alpha_{i-1})$. Hence $A[g^{-1}]$ is a free $B$-module. However, $B$ is not a field; for
$k[a_1, \ldots, a_r]$ is a polynomial ring in $r$ variables over $k$, and hence it contains
infinitely many irreducible polynomials (the proof of this is exactly the
same as Euclid's proof that there exist infinitely many primes). Hence,
there is an irreducible polynomial $h \in k[a_1, \ldots, a_r]$ which does not divide
$g$, and then obviously $h^{-1} \notin k[a_1, \ldots, a_r, g^{-1}]$. Therefore $B$ contains an
ideal $I$ with $I \neq 0, B$, and since $A[g^{-1}]$ is a free module over $B$, $IA[g^{-1}]$
is a proper ideal of $A[g^{-1}]$. Thus $A[g^{-1}]$ is not a field. But if $A$ were a
field then we would have $A[g^{-1}] = A$, and hence $A$ is not a field.
\end{proof}

\begin{theorem}\label{thm:5.3}
Let $k$ be a field, and let $m$ be any maximal ideal of the
polynomial ring $k[X_1, \ldots, X_n]$; then the residue class field $k[X_1, \ldots, X_n]/m$
is algebraic over $k$. Hence $m$ can be generated by $n$ elements, and in
particular if $k$ is algebraically closed then $m$ is of the form $m =
(X_1 - a_1, \ldots, X_n - a_n)$ for $a_i \in k$.
\end{theorem}

\begin{proof}
Set $k[X_1, \ldots, X_n]/m = K$, and write $\alpha_i$ for the image of $X_i$ in $K$; then
$K = k[\alpha_1, \ldots, \alpha_n]$. By the previous theorem, since $K$ is a field it is algebraic
over $k$, and then by Theorem 1, (ii), $m$ is generated by $n$ elements. If $k$ is
algebraically closed then $k = K$, so that each $X_i$ is congruent modulo $m$ to
some $a_i \in k$; then $(X_1 - a_1, \ldots, X_n - a_n) \subset m$. On the other hand
$(X_1 - a_1, \ldots, X_n - a_n)$ is obviously a maximal ideal, so that equality must
hold.
\end{proof}

Let $k$ be a field and $\bar{k}$ its algebraic closure. Suppose that $\Phi \subset
k[X_1, \ldots, X_n]$ is a subset. An $n$-tuple $\alpha = (a_1, \ldots, a_n)$ of elements $a_i \in \bar{k}$ is an
\emph{algebraic zero}\index{algebraic zero} of $\Phi$ if it satisfies $f(\alpha) = 0$ for every $f(X) \in \Phi$.

\begin{theorem}[The Hilbert Nullstellensatz]\label{thm:5.4}
\begin{enumerate}[(i)]
  \item If $\Phi$ is a subset of $k[X_1, \ldots, X_n]$ which does not have any algebraic
        zeros then the ideal generated by $\Phi$ contains $1$.
  \item Given a subset $\Phi$ of $k[X_1, \ldots, X_n]$ and an element $f \in k[X_1, \ldots, X_n]$,
        suppose that $f$ vanishes at every algebraic zero of $\Phi$. Then some power of $f$
        belongs to the ideal generated by $\Phi$, that is there exist $\nu > 0$,
        $g_i \in k[X_1, \ldots, X_n]$ and $h_i \in \Phi$ such that $f^\nu = \sum g_i h_i$.
\end{enumerate}
\end{theorem}

\begin{proof}
(i) Let $I$ be the ideal generated by $\Phi$; if $1 \notin I$ then there exists a
maximal ideal $m$ containing $I$. By the previous theorem, $k[X_1, \ldots, X_n]/m$
is algebraic over $k$, so that it has a $k$-linear isomorphic embedding $\theta$ into
$\bar{k}$. If we set $\theta(X_i \bmod m) = a_i$ then for all $g(X) \in m$ we have $0 = \theta(g(X)) =
g(a)$, therefore $\alpha = (a_1, \ldots, a_n)$ is an algebraic zero of $m$, and
hence also of $\Phi$. This contradicts the hypothesis; hence, $1 \in I$.

(ii) Inside $k[X_1, \ldots, X_n, Y]$ we consider the set $\Phi \cup \{1 - Yf(X)\}$; then this
set has no algebraic zeros, so that by (i) it generates the ideal $(1)$. Therefore
there exists a relation of the form
\[
  1 = \sum p_i(X, Y)h_i(X) + Q(X, Y)(1 - Yf(X)),
\]
with $h_i(X) \in \Phi$. This is an identity in $X_1, \ldots, X_n$ and $Y$, so that it still holds if
we substitute $Y = f(X)^{-1}$. Hence we have
\[
  1 = \sum p_i(X, f^{-1})h_i(X),
\]
so that multiplying through by a suitable power of $f$ and clearing
denominators gives $f^\nu = \sum g_i(X)h_i(X)$, with $g_i \in k[X_1, \ldots, X_n]$ and
$h_i \in \Phi$.
\end{proof}

\begin{remark}
The above proof of (ii) is a classical idea due to Rabinowitch
[1]. In a modern form it can be given as follows: let $I \subset k[X_1, \ldots, X_n] = A$
be the ideal generated by $\Phi$; then in the localisation $A_f$ with respect to $f$ (see
\S\ref{sec:4}, Example 1), we have $IA_f = A_f$, so that a power of $f$ is in $I$.
\end{remark}

\begin{theorem}\label{thm:5.5}
Let $k$ be a field, $A$ a ring which is finitely generated over $k$, and
$I$ a proper ideal of $A$; then the radical of $I$ is the intersection of all maximal
ideals containing $I$, that is $\sqrt{I} = \bigcap_{I \subset m} m$.
\end{theorem}

\begin{proof}
Let $A = k[a_1, \ldots, a_n]$, so that $A$ is a quotient of $k[X_1, \ldots, X_n]$.
Considering the inverse image of $I$ in $k[X]$ reduces to the case $A = k[X]$,
and the assertion follows from Theorem 4, (ii).
\end{proof}

Compared to the result $\sqrt{I} = \bigcap_{I \subset P} P$ proved in \S\ref{sec:1}, the conclusion of
Theorem 5 is much stronger. It is equivalent to the condition on a ring
that every prime ideal $P$ should be expressible as an intersection of maximal
ideals. Rings for which this holds are called \emph{Hilbert rings}\index{Hilbert ring} or \emph{Jacobson rings}\index{Jacobson ring},
and they have been studied independently by O. Goldmann [1] and
W. Krull [7]. See also Kaplansky [K] and Bourbaki [B-5].

\begin{theorem}\label{thm:5.6}
Let $k$ be a field and $A$ an integral domain which is finitely
generated over $k$; then
\[
  \dim A = \trdeg_k A.
\]
\end{theorem}

\begin{proof}
Let $A = k[X_1, \ldots, X_n]/P$, and set $r = \trdeg_k A$. To prove that
$r \ge \dim A$ it is enough to show that if $P$ and $Q$ are prime ideals of
$k[X] = k[X_1, \ldots, X_n]$ with $Q \supset P$ and $Q \neq P$ then
\[
  \trdeg_k k[X]/Q < \trdeg_k k[X]/P.
\]
The $k$-algebra homomorphism $k[X]/P \to k[X]/Q$ is onto, so that $\trdeg_k
k[X]/Q \le \trdeg_k k[X]/P$ is obvious. Suppose that equality holds. Let
$k[X]/P = k[\alpha_1, \ldots, \alpha_n]$ and $k[X]/Q = k[\beta_1, \ldots, \beta_n]$; we can assume that
$\beta_1, \ldots, \beta_r$ is a transcendence basis for $k(\beta)$ over $k$. Then $\alpha_1, \ldots, \alpha_r$ are also
algebraically independent over $k$, so that they form a transcendence basis
for $k(\alpha)$ over $k$. Now set $S = k[X_1, \ldots, X_n] - \{0\}$; $S$ is a multiplicative set in
$k[X]$ with $P \cap S = \emptyset$ and $Q \cap S = \emptyset$. Setting $R = k[X_1, \ldots, X_n]$ and
$K = k(X_1, \ldots, X_n)$, we have $R_S = K[X_{r+1}, \ldots, X_n]$, and
\[
  R_S/PR_S \cong k(\alpha_1, \ldots, \alpha_r)[\alpha_{r+1}, \ldots, \alpha_n],
\]
so that $R_S/PR_S$ is algebraic over $K = k(X_1, \ldots, X_n) \cong k(\alpha_1, \ldots, \alpha_r)$, and
therefore by Theorem 1, $PR_S$ is a maximal ideal of $R_S$; but this contradicts
the assumptions $P \subset Q$ with $P \neq Q$ and $Q \cap S = \emptyset$.

Now let us prove that $r \le \dim A$ by induction on $r$. If $r = 0$ then, by
Theorem 1, $A$ is a field, so $\dim A = 0$ and the assertion holds. Now let $r > 0$,
and suppose that $A = k[\alpha_1, \ldots, \alpha_n]$ with $\alpha_1$ transcendental over $k$; setting
$S = k[X_1] - \{0\}$ and $R = k[X_1, \ldots, X_n]$ we get
\[
  R_S = k(X_1)[X_2, \ldots, X_n] \quad \text{and} \quad R_S/PR_S \cong k(X_1)[\alpha_2, \ldots, \alpha_n].
\]
Hence $R_S/PR_S$ has transcendence degree $r - 1$ over $k(X_1)$, so that by
induction $\dim R_S/PR_S \ge r - 1$. Thus there exists a strictly increasing chain
$PR_S = Q_0 \subsetneq Q_1 \subsetneq \cdots \subsetneq Q_{r-1}$ of prime ideals of $R_S$. If we set $P_i = Q_i \cap R$ then
$P_i$ is a prime ideal of $R$ disjoint from $S$; in particular, the residue class of
$X_1$ in $R/P_{r-1}$ is not algebraic over $k$, and so $\trdeg_k R/P_{r-1} \ge 1$. Then
$P_{r-1}$ is not a maximal ideal of $R$ by Theorem 3, and therefore $R$ has a
maximal ideal $P_r$ strictly bigger than $P_{r-1}$. Hence $\dim A = \coht P \ge r$.
\end{proof}

\begin{corollary*}
If $k$ is a field then $\dim k[X_1, \ldots, X_n] = n$.
\end{corollary*}

We now turn to a different topic, the theorem of Forster and Swan on the
number of generators of a module. Let $A$ be a ring and $M$ a finite $A$-module;
for $p \in \Spec A$, write $\kappa(p)$ for the residue field of $A_p$, and let $\mu(p, M)$ denote the
dimension over $\kappa(p)$ of the vector space $M \otimes \kappa(p) = M_p/pM_p$ (in the usual
sense of linear algebra). This is the cardinality of a minimal basis of the
$A_p$-module $M_p$. Hence, if $p \supset p'$ then $\mu(p, M) \ge \mu(p', M)$.

In 1964 the young function-theorist O. Forster surprised the experts in
algebra by proving the following theorem [1].

\begin{theorem}\label{thm:5.7}
Let $A$ be a Noetherian ring and $M$ a finite $A$-module. Set
\[
  b(M) = \sup\{\mu(p, M) + \coht p \mid p \in \Supp M\};
\]
then $M$ can be generated by at most $b(M)$ elements.
\end{theorem}

This theorem is a very important link between the number of local and
global generators. However, there was room for improvement in the bound
for the number of generators, and in no time R. Swan obtained a better
bound. We will prove Swan's bound. For this we need the concept of \emph{j-Spec
A}\index{j-Spec of a ring} introduced by Swan. This is a space having the same irreducible closed
subsets as m-Spec $A$, but has the advantage of having a `generic point', not
present in m-Spec $A$, for every irreducible closed subset.

A prime ideal which can be expressed as an intersection of maximal
ideals is called a \emph{$j$-prime ideal}\index{j-prime ideal}, and we write j-Spec $A$ for the set of all
$j$-prime ideals. We consider j-Spec $A$ also with its topology as a subspace
of $\Spec A$. Set $M = \text{m-Spec } A$ and $J = \text{j-Spec } A$. If $F$ is a closed subset of
$J$ then there is an ideal $I$ of $A$ such that $F = V(I) \cap J$. One sees easily that
a prime ideal $P$ belongs to $F$ if and only if $P$ can be expressed as an
intersection of elements of $F \cap M = V(I) \cap M$. Hence $F$ is determined by
$F \cap M$, so that there is a natural one-to-one correspondence between
closed subsets of $J$ and of $M$. It follows that if $M$ is Noetherian so is $J$,
and they both have the same combinatorial dimension. Now let $B$ be an
irreducible closed subset of $J$, and let $P$ be the intersection of all the
elements of $B$. If $B = V(I) \cap J$ then $I \subset P$ and we can also write
$B = V(P) \cap J$. If $P$ is not a prime ideal then there exist $f, g \in A$ such that
$f \notin P$, $g \notin P$ and $fg \in P$; but then
\[
  B = (V(P + fA) \cap J) \cup (V(P + gA) \cap J),
\]
and by definition of $P$ there is a $Q \in B$ not containing $f$ and a $Q' \in B$ not
containing $g$, which implies that $B$ is reducible, a contradiction. Therefore
$P$ is a prime ideal. Hence $P \in B$ and $B = V(P) \cap J$. This $P$ is called the
\emph{generic point}\index{generic point} of $B$. Conversely if $P$ is any element of $J$ then $V(P) \cap J$ is
an irreducible closed subset of $J$, and is the closure in $J$ of $\{P\}$. We will
write \emph{$j$-dim $p$} for the combinatorial dimension of $V(p) \cap J$.

For a finite $A$-module $M$ and $p \in J$ we set
\[
  b(p, M) =
  \begin{cases}
    0 & \text{if } M_p = 0 \\
    \text{$j$-dim } p + \mu(p, M) & \text{if } M_p \neq 0.
  \end{cases}
\]

\begin{theorem}[Swan~\cite{Sw1}]\label{thm:5.8}
Let $A$ be a ring, and suppose that m-Spec $A$ is
a Noetherian space. Let $M$ be a finite $A$-module. If
\[
  \sup\{b(p, M) \mid p \in \text{j-Spec } A\} = r < \infty
\]
then $M$ is generated by at most $r$ elements.
\end{theorem}

\begin{proof}
\textbf{Step 1.} For $p \in \Spec A$ and $x \in M$, we will say that $x$ is \emph{basic} at $p$ if $x$
has non-zero image in $M \otimes \kappa(p)$. It is easy to see that this condition is
equivalent to $\mu(p, M/Ax) = \mu(p, M) - 1$.

\begin{lemma*}
Let $M$ be a finite $A$-module and $p_1, \ldots, p_n \in \Supp(M)$. Then there
exists $x \in M$ which is basic at each of $p_1, \ldots, p_n$.
\end{lemma*}

\begin{proof}[Proof of the lemma]
By reordering $p_1, \ldots, p_n$ we assume that $p_i$ is maximal among
$\{p_i, p_{i+1}, \ldots, p_n\}$ for each $i$. By induction on $n$ suppose that $x' \in M$ is basic at
$p_1, \ldots, p_{n-1}$. If $x'$ is basic at $p_n$, then we can take $x = x'$. Suppose then that $x'$
is not basic at $p_n$. By assumption $M_{p_n} \neq 0$ so that we can choose some $y \in M$
which is basic at $p_n$. We have $p_1 \cdots p_{n-1} \not\subset p_n$, so that if we take an element
$a \in p_1 \cdots p_{n-1}$ not belonging to $p_n$, and set $x = x' + ay$, this $x$ satisfies our
requirements. This proves the lemma.
\end{proof}

\textbf{Step 2.} Setting $\sup\{b(p, M) \mid p \in \text{j-Spec } A\} = r$, we now show that there are
just a finite number of primes $p$ such that $b(p, M) = r$. Indeed, for
$n = 1, 2, \ldots$, the subset $X_n = \{p \in \text{j-Spec } A \mid \mu(p, M) \ge n\}$ is closed in j-Spec $A$ by
Theorem~\ref{thm:4.10}; it has a finite number of irreducible components (by
Ex.~\ref{ex:4.11-ex}), and we let $p_{ni}$ (for $1 \le i \le v_n$) be their generic points. If $M$ is
generated by $s$ elements then $X_n = \emptyset$ for $n > s$, so that the set $\{p_{ni}\}_{n,i}$ is
finite. Let us prove that if $b(p, M) = r$ then $p \in \{p_{ni}\}_{n,i}$. Suppose
$\mu(p, M) = n$; then $p \in X_n$, so that by construction $p \supset p_{ni}$ for some $i$. But
if $p \neq p_{ni}$ then j-dim $p$ < j-dim $p_{ni}$, and since $\mu(p, M) = n = \mu(p_{ni}, M)$ we have
$b(p, M) < b(p_{ni}, M)$, which is a contradiction. Hence $p = p_{ni}$.

\textbf{Step 3.} Let us choose an element $x \in M$ which is basic for each of the
finitely many primes $p$ with $b(p, M) = r$, and set $\bar{M} = M/Ax$; then clearly
$b(p, \bar{M}) \le r - 1$ for every $p \in \text{j-Spec } A$. Hence by induction $\bar{M}$ is generated by
$r - 1$ elements, and therefore $M$ by $r$ elements.
\end{proof}

Swan's paper contains a proof of the following generalisation to non-
commutative rings: let $A$ be a commutative ring, $\Lambda$ a possibly non-
commutative $A$-algebra and $M$ a finite left $\Lambda$-module. Suppose that
m-Spec $A$ is Noetherian, and that for every maximal ideal $p$ of $A$ the
$\Lambda_p$-module $M_p$ is generated by at most $r$ elements; then $M$ is generated as a
$\Lambda$-module by at most $r + d$ elements, where $d$ is the combinatorial
dimension of m-Spec $A$.

The Forster--Swan theorem is a statement that local properties imply
global ones; remarkable results in this direction have been obtained by
Mohan Kumar [2] (see also Cowsik--Nori [1] and Eisenbud--Evans [1],
[2]). The number of generators of ideals in local rings is the subject of
a nice book by J. Sally [Sa].

% ============ 习题 ============
\section*{Exercises to \S\ref{sec:5}}
\emph{Prove the following propositions.}

\begin{exercise}\label{ex:5.1-ex}
Let $k$ be a field, $R = k[X_1, \ldots, X_n]$ and let $P \in \Spec R$; then $\htop P +
\coht P = n$.
\end{exercise}

\begin{exercise}\label{ex:5.2-ex}
A zero-dimensional Noetherian ring is Artinian (the converse to Example
2 above).
\end{exercise}
